\documentclass{ximera}
%geschreven 9/11/26
%verbeterd 12/11/26
\title{Het magnetisch veld}
\author{Daan Lipkens}

\begin{document}
\begin{abstract}
    Inleiding tot permanente magneten, het magnetisch veld en de magnetische velden van een rechte stroomvoerende geleider en een stroomvoerende spoel.
\end{abstract}
\maketitle
\label{hoofdstuk:magnetisme}
% ----------------------------------------------------------------------
%          inleiding hoofdstuk
% ----------------------------------------------------------------------
\section*{Inleiding}
Een koelkastmagneetje, de naald van een kompas, de motor van een elektrische fiets en een MRI-scanner in het ziekenhuis: overal speelt magnetisme een rol. Net als de elektrische kracht is de magnetische kracht een kracht die op afstand werkt. Een magneet trekt een paperclip aan zonder hem aan te raken.

In de vorige hoofdstukken beschreven we de elektrische krachtwerking met een \emph{elektrisch veld} rond ladingen. In dit hoofdstuk doen we hetzelfde voor magnetisme. We starten bij de permanente magneten en hun eigenschappen, zie~\verwijs[Permanente magneten]{sec:permanente-magneten}{magnetisme_permanente_magneten}. Daarna voeren we het \emph{magnetisch veld} in en stellen we het voor met veldlijnen, zie~\verwijs[Het magnetisch veld en veldlijnen]{sec:magnetisch-veld}{magnetisme_veld}.

In 1820 ontdekte Hans Christian Oersted dat een kompasnaald uitwijkt in de buurt van een draad waardoor een stroom loopt. Een elektrische stroom maakt dus ook een magnetisch veld. We bekijken het veld van een rechte stroomvoerende geleider, zie~\verwijs[Het magnetisch veld van een rechte stroomvoerende geleider]{sec:veld-rechte-geleider}{magnetisme_rechte_geleider}, en van een stroomvoerende spoel, zie~\verwijs[Het magnetisch veld van een stroomvoerende spoel]{sec:veld-spoel}{magnetisme_spoel}. We sluiten af met de elektromagneet, een spoel met een ijzeren kern die je kan aan- en uitschakelen, zie~\verwijs[Elektromagneten]{sec:elektromagneet}{magnetisme_elektromagneet}.

Doorheen het hoofdstuk zal blijken dat er een sterk verband is tussen elektriciteit en magnetisme: \textbf{bewegende ladingen zijn de bron van elk magnetisch veld}.

\begin{definition}[title={Wat je na dit hoofdstuk moet kennen en kunnen}]\nl
    \begin{itemize}
        \item de eigenschappen van permanente magneten opsommen en magnetische influentie uitleggen;
        \item aangeven welke stoffen door een magneet aangetrokken worden (ferromagnetische stoffen);
        \item het magnetisch veld in een punt beschrijven met een kompasnaald (richting, zin en grootte in tesla);
        \item de veldlijnenpatronen tekenen en interpreteren bij permanente magneten, bij een rechte stroomvoerende geleider en bij een stroomvoerende spoel;
        \item de zin van de veldlijnen bij een rechte geleider en bij een spoel bepalen met de rechterhandregels;
        \item de grootte van het magnetisch veld berekenen bij een rechte stroomvoerende geleider en in een stroomvoerende spoel;
        \item het resulterende magnetisch veld bepalen als de velden op één lijn liggen of loodrecht op elkaar staan;
        \item de analogie en de verschillen geven tussen het zwaarteveld, het elektrisch veld en het magnetisch veld;
        \item uitleggen waarom we besluiten dat bewegende ladingen de bron zijn van magnetisme;
        \item de werking en de voordelen van een elektromagneet uitleggen.
    \end{itemize}
\end{definition}

\end{document}
